\section{アーキテクチャ}

\subsection{全体構成}

本ツールは純粋なクライアントサイド Single Page Application (SPA) であり、
サーバー通信は行わない（フォントファイル取得を除く）。

\begin{center}
\begin{tikzpicture}[
  node distance=3mm and 4mm,
  box/.style={draw=accentblue,fill=accentblue!8,rounded corners=1.5mm,
              minimum width=34mm,minimum height=10mm,align=center,
              font=\small},
  core/.style={draw=headerblue,fill=headerblue!10,rounded corners=1.5mm,
               minimum width=34mm,minimum height=10mm,align=center,
               font=\small\bfseries},
  arr/.style={-{Latex[length=2mm]},draw=headerblue,thick},
]
  \node[box] (user)   {ユーザー操作\\(UI)};
  \node[box,right=of user] (state)  {useAppState\\(useReducer)};
  \node[core,right=of state] (core) {core層\\(rasterize/encode/format)};
  \node[box,below=of user]  (persist) {usePersist\\(LocalStorage)};
  \node[box,below=of state] (share)   {share.ts\\(URL クエリ)};
  \node[box,below=of core]  (canvas)  {Canvas API\\FontFace API};
  \node[box,right=of core]  (output)  {OutputPanel\\(クリップボード)};

  \draw[arr] (user) -- (state);
  \draw[arr] (state) -- (core);
  \draw[arr] (core) -- (output);
  \draw[arr] (state) -- (persist);
  \draw[arr] (state) -- (share);
  \draw[arr] (core) -- (canvas);
\end{tikzpicture}
\end{center}

\subsection{レイヤ分割}

\begin{tabularx}{\textwidth}{l X}
\toprule
\textbf{レイヤ} & \textbf{責務} \\
\midrule
\texttt{src/core/} & 純粋関数群。DOM依存を最小化（Canvas利用箇所のみ）。ラスタライズ、ビットパック、フォーマッタ、共有URLエンコード。\\
\texttt{src/hooks/} & React フック。状態管理（useReducer）、永続化（LocalStorage）、テーマ切替、グリフメモ化。\\
\texttt{src/components/} & UI コンポーネント。shadcn/ui プリミティブを基盤とし、パネル単位で責務分離。\\
\texttt{src/components/ui/} & Radix UI ベースの汎用プリミティブ（Button, Card, Input, Tabs, Switch, Slider, Select, Tooltip 等）。\\
\bottomrule
\end{tabularx}

\subsection{データフロー}

\begin{enumerate}[leftmargin=*,itemsep=2pt]
  \item ユーザー入力 → \texttt{useAppState} の dispatch でアクション発行。
  \item Reducer が新しい \texttt{AppState} を返し、コンポーネントが再描画。
  \item \texttt{useGlyphs} が入力テキストとフォント設定に基づき
        \texttt{rasterizeChar()} を再実行しグリフ配列をメモ化。
  \item \texttt{formatOutput()} がグリフ配列 + フォーマット設定から
        出力文字列を生成し、\texttt{OutputPanel} が表示。
  \item 同時に \texttt{usePersist} が \texttt{AppState} を
        LocalStorage に書き戻す（\texttt{overrides} を除く）。
\end{enumerate}

\subsection{ビルド・配布パイプライン}

\begin{enumerate}[leftmargin=*,itemsep=2pt]
  \item 開発: \texttt{npm run dev}（Vite 開発サーバ、HMR有効）。
  \item 本番ビルド: \texttt{npm run build}（Vite が \texttt{dist/} を生成）。
  \item 自動デプロイ: \texttt{main} ブランチ push を契機に
        \texttt{.github/workflows/deploy.yml} が起動、
        GitHub Pages（\texttt{gh-pages} 相当の Actions 経由）へ配信。
  \item SPA フォールバック: デプロイ時に \texttt{404.html} を
        \texttt{index.html} と同内容でコピーし、
        任意パスでのアクセスに対応。
\end{enumerate}

\newpage
